% from ega0-21-fr.tex
\oldpage[0\textsubscript{IV-1}]{176}
Il suffit de prouver qu'un élément $f\ne0$ du premier membre de (21.8.8.4)
appartient au second membre. Soient $\gamma$ un indice,
$g\in k_\gamma[[T_1,\ldots,T_r]]$ tel que
$gf\in k_\gamma[[T_1,\ldots,T_r]]$~; quitte à remplacer $g$ par $g^q$, on
peut supposer que $g\in k_0[[T_1,\ldots,T_r]]$. Alors, pour tout
$\alpha\geqslant\gamma$, on a
\[
  gf\in k_\gamma[[T_1,\ldots,T_r]]\cap k_\alpha((T_1,\ldots,T_r))
  =k_\alpha[[T_1,\ldots,T_r]]
\]
en vertu du lemme (21.8.8.2). Mais il est clair que l'intersection des
anneaux $k_\alpha[[T_1,\ldots,T_r]]$ n'est autre que
$k_0[[T_1,\ldots,T_r]]$, et l'on a donc bien
$f\in k_0((T_1,\ldots,T_r))$.
\end{proof}

\subsection{Modules de différentielles complétés dans les anneaux de séries formelles}
\label{subsection:0.21.9-fr}

Dans ce numéro, les corps ne sont plus nécessairement supposés être de
caractéristique $>0$.

\begin{lemma}[21.9.1]
\label{0.21.9.1-fr}
Soient $k$ un corps, $A$ un anneau local noethérien complet qui est une
$k$-algèbre, $K$ le corps résiduel de $A$. Si $K$ est une extension de type
fini de $k$, le $A$-module $\widehat\Omega^1_{A/k}$ est de type fini.
\end{lemma}

\begin{proof}
En vertu de (20.7.15), il suffit de prouver que $\Omega^1_{K/k}$ est un
$K$-espace vectoriel de rang fini. Or, par hypothèse, $K$ est le corps des
fractions d'une $k$-algèbre de type fini $B$~; comme $\Omega^1_{B/k}$ est un
$B$-module de type fini en vertu de (20.4.7), la conclusion résulte de
(20.5.9).
\end{proof}

\begin{proposition}[21.9.2]
\label{0.21.9.2-fr}
Soient $k_0$ un corps, $A$ un anneau local noethérien complet qui est une
$k_0$-algèbre formellement lisse, $\mathfrak m$ son idéal maximal,
$K=A/\mathfrak m$ son corps résiduel (de sorte qu'en tant qu'anneau $A$ est
isomorphe à un anneau de séries formelles $K[[T_1,\ldots,T_r]]$ (19.6.5)).
Si $K$ est une extension de type fini de $k_0$, alors
$\widehat\Omega^1_{A/k_0}$ est un $A$-module libre de rang égal à
$\dim(A)+\deg.\operatorname{tr}_{k_0}K$. En outre, pour tout sous-corps $k$
de $k_0$ tel que $\Omega^1_{k_0/k}$ soit de rang fini,
$\widehat\Omega^1_{A/k}$ est un $A$-module libre de rang égal à
\[
  \dim(A)+\deg.\operatorname{tr}_{k_0}K
  +\operatorname{rg}_{k_0}(\Omega^1_{k_0/k}).
\]
\end{proposition}

\begin{proof}
Il est clair que, si $P$ est le sous-corps premier de $k_0$, $K$ est une
$P$-algèbre formellement lisse (19.6.1). Notons d'autre part que l'on a
$\Upsilon^K_{A/k_0/P}=0$~: en effet, cela revient à voir que l'homomorphisme
\[
  u_{A/k_0/P}\otimes1_K:
  \Omega^1_{k_0/P}\otimes_{k_0}K
  \longrightarrow \Omega^1_{A/P}\otimes_A K
\]
est injectif (20.6.14.5), $u$ désignant l'homomorphisme structural
$k_0\to A$. Mais comme $A$ est une $k_0$-algèbre formellement lisse par
hypothèse, il résulte de (20.7.2) que $u_{A/k_0/P}$ est formellement
inversible à gauche, donc (19.1.7) $u_{A/k_0/P}\otimes1_K$ est inversible à
gauche et \emph{a fortiori} injectif, ce qui prouve notre assertion.
Appliquant la suite exacte (20.6.22.1), où $\Lambda,A,B,C$ sont remplacés par
$P,k_0,A,K$, il vient la suite exacte
\[
  0\longrightarrow \Upsilon^K_{A/k_0}
  \longrightarrow \mathfrak m/\mathfrak m^2
  \longrightarrow \Omega^1_{A/k_0}\otimes_A K
  \longrightarrow \Omega^1_{K/k_0}\longrightarrow0,
\]
de sorte que l'on a
\[
\begin{split}
  \operatorname{rg}_K(\Omega^1_{A/k_0}\otimes_AK)
  &={\operatorname{rg}}_K(\mathfrak m/\mathfrak m^2)
    +\bigl(\operatorname{rg}_K(\Omega^1_{K/k_0})
    -\operatorname{rg}_K(\Upsilon_{K/k_0})\bigr)\\
  &={\operatorname{rg}}_K(\mathfrak m/\mathfrak m^2)
    +\deg.\operatorname{tr}_{k_0}K
\end{split}
\]
en vertu de (21.7.1). D'autre part, $\widehat\Omega^1_{A/k_0}$ est un
$A$-module formellement projectif (20.4.9), et comme sa topologie est la
topologie $\mathfrak m$-préadique (20.4.5),
$\widehat\Omega^1_{A/k_0}/\mathfrak m^{j+1}\widehat\Omega^1_{A/k_0}$ est un
$(A/\mathfrak m^{j+1})$-module projectif pour tout $j$ (19.2.4), donc libre
($\mathbf 0_{\mathrm{III}}$, 10.1.2) et de rang

